% Source: https://artofproblemsolving.com/community/c3h1375394_geo_questions
% Stable identity: aops-topic-1375394-q17
Place the base in the plane $z=0$ with $A=(0,0,0)$, $B=(12,0,0)$, $C=(12,12,0)$, $D=(0,12,0)$, and let $E=(x,y,h)$ with $h>0$. Let $H=(x,y,0)$ be the foot of the height.

\begin{asy}
usepackage("amsmath");
size(6.5cm);
defaultpen(fontsize(10pt));
real h=3*sqrt(7);
pair pr(real x,real y,real z){return (x-0.3*y,0.3*y+z);}
pair A=pr(0,0,0),B=pr(12,0,0),C=pr(12,12,0),D=pr(0,12,0),T=pr(9,7,h),H=pr(9,7,0);
draw(D--A--B,linewidth(0.9)); draw(B--C--D,dashed);
draw(T--A,linewidth(0.9)); draw(T--B,linewidth(0.9)); draw(T--D,linewidth(0.9)); draw(T--C,dashed);
draw(T--H,red+dashed); dot(H,red);
label("$A$",A,SW); label("$B$",B,SE); label("$C$",C,dir(240)); label("$D$",D,W); label("$E$",T,N);
label("$H$",H,SW,red); label("$h$",H+(0,2.6),W,red);
label("$12$",(A+B)/2,S); label("$12$",(A+D)/2,SW);
label("$\sqrt{193}$",A+0.22*(T-A)+(-0.35,0),W,fontsize(9pt)); label("$11$",B+0.5*(T-B),E,fontsize(9pt));
label("$13$",D+0.5*(T-D),N,fontsize(9pt));
label("\mbox{}",(min(currentpicture,user=true)+max(currentpicture,user=true))/2,black); // keeps the text after the figure black when LaTeX typesets the labels inline
\end{asy}

Then
\[
AE^2=x^2+y^2+h^2=193,\qquad BE^2=(x-12)^2+y^2+h^2=121,\qquad DE^2=x^2+(y-12)^2+h^2=169 .
\]
Subtracting, $AE^2-BE^2=24x-144=72$, so $x=9$, and $AE^2-DE^2=24y-144=24$, so $y=7$. Then $h^2=193-81-49=63$ and $h=3\sqrt7$.

The fourth length is consistent: $CE^2=(9-12)^2+(7-12)^2+63=9+25+63=97$. (In general $AE^2+CE^2=BE^2+DE^2$ for any point $E$ and rectangle $ABCD$; here $193+97=121+169=290$.)

The volume is
\[
V=\dfrac{1}{3}\cdot12^2\cdot3\sqrt7=\boxed{144\sqrt7}.
\]
